\section{Introduction}
\label{sec:intro}

Recently, the rapid acceleration of large language model (LLM) development has transformed artificial intelligence (AI) from a narrow, task-specific paradigm into a general-purpose intelligence with far-reaching implications. Contemporary LLMs demonstrate increasingly sophisticated linguistic, reasoning, and problem-solving capabilities, and are showcasing superb human-like behaviors, prompting a fundamental research question: 

\begin{quote}
\begin{center}
\emph{To what extent do these systems exhibit forms of awareness?}
\end{center}
\end{quote}

We want to note here, while the notion of AI consciousness remains both philosophically contentious and empirically less grounded, the concept of AI awareness --- defined as a system’s functional ability to represent and reason about its own identity, capabilities, and informational states --- has emerged as an important and tractable research frontier.

The theoretical foundations of awareness are deeply rooted in cognitive science and psychology, where awareness is typically classified into perceptual awareness, self-awareness, and social awareness. In these fields, awareness is understood as the ability to access and monitor mental states, reason about them, and adjust behavior accordingly. Recent work in computational cognitive science suggests that certain aspects of awareness may be approximated by large-scale generative models that exhibit metacognitive behaviors, calibrate their epistemic confidence, and engage in perspective-taking. These functional abilities raise pressing questions about how awareness manifests in LLMs, how it can be systematically measured, and what implications it holds for AI capability, safety and alignment.

Despite the increasing attention devoted to this topic, the study of AI awareness remains fragmented across disciplines, with limited consensus regarding its definition, measurement, and significance. Some scholars emphasize emergent capabilities observed through prompt-based introspection or theory-of-mind-inspired tests; others caution against anthropomorphizing statistical models, warning that apparent self-reflection may arise from linguistic pattern completion rather than genuine metacognitive representation. Moreover, current methodologies for assessing AI awareness are often limited by confounding factors such as prompt sensitivity, dataset leakage, and lack of longitudinal robustness.


This review provides the first comprehensive, cross-disciplinary synthesis of AI awareness research. We begin by laying out the theoretical foundations, carefully distinguishing AI awareness from AI consciousness, and examining how awareness-related concepts have been formalized in cognitive and computational sciences. Next, we critically evaluate experimental approaches for testing LLM awareness, highlighting both their empirical findings and methodological limitations. We then explore the potential beneficial relationships between functional AI awareness and AI capabilities, from improving model reasoning and task planning to enhancing AI safety. Finally, we address the emerging risks associated with more aware AI systems, focusing on key concerns within the safety and alignment community---such as risks of deception, manipulation, and emergent behaviors that challenge controllability---and broader ethical issues, including concerns like false anthropomorphization.



By integrating insights from AI, cognitive science, psychology, AI safety and alignment, this review not only provides a structured and authoritative account of the current state of knowledge but also charts a roadmap for future inquiry. In doing so, it seeks to advance understanding of one of the most profound and interdisciplinary challenges at the intersection of AI, cognition science and our society.

